\begin{onehalfspace}

Questa sezione raccoglie il materiale di dettaglio del Capitolo~\ref{chap:obiettivi}: motivazione crittografica, formulazione analitica degli obiettivi, contributi, requisiti, use case, input/output, metriche e test. Ogni blocco dichiara la propria origine nel corpo e completa la sintesi senza modificarne il perimetro.

\subsection{Motivazione: la Vulnerabilità della Crittografia Attuale}
\label{sec:motivazione}

La quasi totalità delle comunicazioni digitali sicure --- dalle transazioni bancarie alle connessioni \ac{HTTPS}, dai protocolli \ac{SSH} ai certificati digitali --- è protetta da sistemi crittografici la cui sicurezza si fonda su un presupposto computazionale: la difficoltà di fattorizzare il prodotto di due grandi numeri primi. Il sistema \ac{RSA}, ad esempio, costruisce la propria chiave pubblica come $N = p \cdot q$, dove $p$ e $q$ sono primi di centinaia o migliaia di bit. Finché non esiste un algoritmo classico efficiente per recuperare $p$ e $q$ a partire da $N$, la chiave privata rimane al sicuro.

Per comprendere perché i sistemi crittografici attuali siano strutturalmente vulnerabili all'algoritmo di Shor, è necessario mettere a confronto le complessità computazionali in gioco. Il miglior algoritmo classico per la fattorizzazione, il \ac{GNFS}, richiede un tempo:
\begin{equation}
    T_{\text{classico}}(N) = O\!\left(\exp\!\left((c+o(1))(\log N)^{1/3}(\log\log N)^{2/3}\right)\right)
    \label{eq:gnfs_obiettivi}
\end{equation}
che, pur crescendo più lentamente di un esponenziale puro, rimane del tutto impraticabile per i valori di $N$ usati in crittografia. A titolo di esempio, fattorizzare un numero RSA a 2048 bit con algoritmi classici richiederebbe un tempo stimato superiore all'età dell'universo, anche impiegando l'intera potenza di calcolo disponibile sul pianeta.

Nel 1994, Peter Shor dimostrò che un computer quantistico è in grado di eseguire la stessa operazione in tempo \textbf{polinomiale}~\cite{shor1994}:
\begin{equation}
    T_{\text{Shor}}(N) = O\!\left((\log N)^3\right)
    \label{eq:shor_obiettivi}
\end{equation}
Lo scarto tra le due espressioni non è una differenza di grado, ma una differenza di \textit{natura}: il tempo classico cresce in modo super-polinomiale con la lunghezza della chiave, mentre il tempo quantistico cresce polinomialmente. \textbf{Aumentare la lunghezza della chiave RSA non è quindi una contromisura strutturale} contro Shor: raddoppiare $n$ porta il costo asintotico textbook da $O(n^3)$ a $8\,O(n^3)$. Sul piano concreto questo fattore può richiedere molte più risorse fault-tolerant e ritardare un attacco, ma non ripristina un problema computazionalmente difficile per un avversario quantistico scalabile. La stessa vulnerabilità asintotica si estende ai sistemi basati sul logaritmo discreto (Diffie--Hellman, \ac{ECDSA}).

La minaccia ha già una conseguenza operativa: nel modello \textit{harvest now, decrypt later}, dati cifrati oggi e destinati a restare sensibili vengono conservati in vista di una futura macchina crittograficamente rilevante. Non esiste però una previsione affidabile della sua data di arrivo. Le roadmap dei produttori descrivono obiettivi ingegneristici e il \ac{NIST} sottolinea che le stime spaziano da pochi anni a decenni; proprio l'incertezza, unita ai tempi lunghi di migrazione, motiva l'adozione anticipata della \ac{PQC}~\cite{nist_pqc_explainer2026,ibm_roadmap2025}.

Gli ostacoli pratici sono la scala, la profondità fault-tolerant e il \textbf{rumore}. I processori attuali operano in regime \ac{NISQ}: dispongono di molti qubit fisici, ma decoerenza, errori di gate, readout, connettività e overhead di correzione impediscono di sostenere il numero di operazioni logiche richiesto da Shor su chiavi crittografiche. Comprendere questi limiti e stabilire quali strategie li riducano --- e in quale punto della pipeline i modelli di machine learning contribuiscano davvero --- è la ragione d'essere del presente lavoro.

\subsection{Approfondimento della Sezione 2.2: Obiettivi Specifici}
\label{app:obiettivi:specifici}

\noindent\textit{Origine nel corpo:} formulazione estesa degli obiettivi sintetizzati nella Sezione ``Obiettivi specifici'' del Capitolo~\ref{chap:obiettivi}.

\begin{enumerate}
    \item \textbf{Fondamenti del calcolo quantistico.} Fornire una trattazione rigorosa del qubit, della sovrapposizione, dell'entanglement, delle porte logiche e dell'interferenza, base necessaria per comprendere l'algoritmo discusso nella Sezione~\ref{chap:shor}.

    \item \textbf{Analisi teorica dell'algoritmo di Shor.} Descrivere la riduzione della fattorizzazione al \textit{period finding}, la \ac{QFT} e la phase estimation; analizzare la complessità polinomiale e le conseguenze per \ac{RSA} e la sicurezza post-quantum.

    \item \textbf{Implementazione su Qiskit e analisi sperimentale.} Costruire la \ac{QFT}, l'esponenziazione modulare e il circuito completo, verificarne il comportamento ideale e analizzarne le prestazioni sulle istanze dichiarate nel perimetro.

    \item \textbf{Rumore e strategie di protezione.} Studiare decoerenza, errori di gate e readout, con attenzione a \ac{QFT} e phase estimation; distinguere e valutare \ac{QEC}, \ac{ZNE}, \ac{PEC} e mitigazione del readout, senza attribuire a una tecnica risultati non prodotti dagli esperimenti.

    \item \textbf{Ottimizzazione mediante modelli appresi.} Stabilire in quale punto della pipeline un modello migliori effettivamente la resa. A valle dell'esecuzione, il classificatore sull'istogramma è sottoposto ad ablazione contro TOP-K; durante la correzione, il decoder addestrato sulle sindromi è confrontato con il riferimento analitico a parità di campioni. Il criterio è quantitativo --- iterazioni nella prima fase, logical error rate nella seconda --- e mira a delimitare sia il regime di utilità sia quello di non utilità.

    \item \textbf{Correzione degli errori e integrazione con Shor.} Costruire in ordine di complessità codice a ripetizione, Steane $[[7,1,3]]$ e surface code; misurare la corrispondenza sindrome--errore, la soppressione di $p_L$ in funzione di $p$ e $d$ e il comportamento di soglia. Studiare poi Shor sotto un residuo Pauli fenomenologico per gate, senza identificarlo direttamente con le metriche aggregate dei benchmark QEC.
\end{enumerate}

Questi obiettivi producono tre classi di output: una ricostruzione teorica verificabile, una piattaforma sperimentale riproducibile e confronti quantitativi contro baseline esplicite. Le domande DR1--DR6 e le ipotesi restano nel corpo, dove costituiscono l'indice logico dei risultati.

\subsection{Approfondimento della Sezione 2.4: Contributi e Confini Interpretativi}
\label{app:obiettivi:contributi}

\noindent\textit{Origine nel corpo:} dettaglio dei quattro contributi riassunti nella Sezione~\ref{sec:contributo_originale}.

\textbf{Progettazione e verifica di TOP-K.} Rispetto alle tecniche strutturali applicate prima della misura, la prima fase valuta una regola di post-processing direttamente sull'istogramma rumoroso. Il protocollo appaiato separa il contributo della ricerca multi-candidato da quello del classificatore senza confonderlo con una diversa estrazione Monte Carlo.

\textbf{Quantificazione del fattore $\rho$.} I due preset illustrativi usano lo stesso circuito $N=15$; per ogni coppia replica--iterazione, M1, TOP-4 e M2 ricevono lo stesso istogramma. Le differenze sono analizzate con Wilcoxon appaiato--Pratt. Le istanze Beauregard più grandi forniscono controlli ideali di correttezza e costo, non osservazioni delle campagne rumorose. I risultati quantitativi provengono soltanto dagli artefatti schema~2 generati dopo la correzione del moltiplicatore $\times7\bmod15$.

\textbf{Caratterizzazione a livelli.} Ripetizione, Steane e surface code misurano il logical error rate nei rispettivi protocolli; M8 misura separatamente la sensibilità di Shor a un proxy Pauli uniforme per gate. In assenza di una compilazione fault-tolerant, la tesi non sovrappone sulla curva M8 tassi per ciclo o per memoria e non confronta codici sotto unità non commensurabili.

\textbf{Regime di utilità del machine learning.} A valle di Shor, M1--TOP-4--M2 separa il classificatore dalla ricerca multi-candidato. Nel surface code, modello appreso e decoder analitico ricevono gli stessi detection events. Questa simmetria sperimentale permette di esprimere il criterio di applicabilità in termini verificabili anziché attribuire il risultato alla sola scelta dell'algoritmo di apprendimento.

\subsection{Requisiti Funzionali}
\label{sec:requisiti}

Definiti gli obiettivi e il contributo originale, è possibile formalizzare cosa il sistema sperimentale deve concretamente fare. I requisiti seguenti traducono le ipotesi di ricerca in vincoli tecnici verificabili, e costituiscono il riferimento rispetto al quale l'implementazione descritta nel Capitolo~\ref{chap:sviluppo} viene valutata nel Capitolo~\ref{chap:risultati}.

Il sistema sperimentale deve soddisfare i seguenti requisiti funzionali:

\begin{enumerate}

    \item \textbf{Esecuzione configurabile entro il perimetro validato.}
    La campagna rumorosa e la demo pubblica usano l'istanza canonica $N=15$, $a=7$ e un numero configurabile di shot entro limiti dichiarati. Le implementazioni Beauregard per $N=21$ e $N=35$ sono sottoposte a test ideali e audit di risorse, ma non sono esposte come campagne rumorose equivalenti.

    \item \textbf{Modello di rumore parametrizzabile.}
    Il sistema permette di configurare separatamente i parametri Aer $\lambda_{1q}$ e $\lambda_{2q}$, i tempi $T_1/T_2$, il readout asimmetrico della demo e l'eventuale sovrarotazione coerente. I preset della campagna restano uniformi e illustrativi; non riproducono una calibrazione hardware nominata.

    \item \textbf{Raccolta e archiviazione degli output.}
    Il sistema registra la distribuzione di misura del registro di controllo per ogni esecuzione e la archivia in formato strutturato. La riproducibilità usa un seed base e calendari derivati separati per replica, iterazione, simulatore e risoluzione dei pareggi.

    \item \textbf{Analisi statistica del Metodo 1.}
    Il sistema implementa la pipeline del Metodo 1: selezione della misura più frequente dell'istogramma (strategia TOP-1, senza sogliatura preliminare), stima del periodo $r$ tramite frazioni continue, verifica della correttezza dei fattori estratti. Misura e registra il numero di iterazioni necessarie per raggiungere il primo risultato corretto.

    \item \textbf{Training, validazione e inferenza del Metodo 2.}
    Il sistema genera il dataset dal simulatore rumoroso, addestra il classificatore con divisione training/validation/test e ne valuta le metriche. Il modello corrente apprende l'etichetta TOP-1: decide se il candidato dominante conduce ai fattori; se accetta l'istogramma, M2 applica poi la ricerca TOP-4. L'audit TOP-16 è a classe singola e non produce un classificatore.

    \item \textbf{Confronto sistematico sui use case previsti.}
    M1, TOP-4 e M2 vengono eseguiti sugli stessi istogrammi nei due preset $N=15$, producendo metriche appaiate secondo la Sezione~\ref{sec:metriche}. Per $N=21$ e $N=35$ si riportano correttezza ideale e costo compilato; non si afferma che la probabilità rumorosa sia zero, perché la campagna rumorosa non viene eseguita.

    \item \textbf{Riproducibilità.}
    Tutti gli esperimenti sono riproducibili: seed fisso, parametri esplicitati, codice disponibile nel repository del progetto. Questo requisito è condizione necessaria per la validazione accademica dei risultati \cite{shor_hpec2023}.

\end{enumerate}

\subsection{Approfondimento della Sezione 2.5: Campagne e Validazioni}
\label{app:obiettivi:usecase}

\noindent\textit{Origine nel corpo:} specifica completa del perimetro sintetizzato nella Sezione~\ref{sec:use_case}.

Il disegno separa le campagne rumorose dalle validazioni ideali:

\begin{itemize}
    \item \textbf{Campagne rumorose N15}: stesso $N=15$, stessa base $a=7$, stesso circuito e due preset, \emph{riferimento} e \emph{stress}. Sono scenari uniformi illustrativi, non calibrazioni di una \ac{QPU} reale.
    \item \textbf{Validazioni ideali N21/N35}: i moltiplicatori modulari di Beauregard sono verificati senza rumore mediante truth table, ripristino dei registri ausiliari e prove QPE ideali. Il loro costo esclude una campagna Aer rumorosa equivalente; l'assenza di tale campagna non viene interpretata come probabilità rumorosa nulla.
\end{itemize}

\subsubsection{Preset delle campagne rumorose}

La Tabella~\ref{tab:use_case_params} riporta le sole configurazioni rumorose. I simboli $\epsilon_{1q}$ e $\epsilon_{2q}$ indicano esattamente il parametro $\lambda$ passato da Aer a \texttt{depolarizing\_error}; non sono probabilità Pauli aggregate né dati di calibrazione hardware. Il modello è uniforme su tutti i qubit.

\begin{table}[H]
\centering
\resizebox{\textwidth}{!}{%
\begin{tabular}{@{}lcc@{}}
\toprule
\textbf{Parametro} & \textbf{Riferimento} & \textbf{Stress} \\
\midrule
$N$                          & 15  & 15  \\
$a$                          & 7   & 7   \\
Periodo $r$ atteso           & 4   & 4   \\
$n_{\text{count}}$ (qubit)  & 8   & 8   \\
$M_{\text{shots}}$          & 1024 & 1024 \\
\midrule
$\epsilon_{1q}=\lambda_{1q}$ & $1\times10^{-3}$ & $5\times10^{-3}$ \\
$\epsilon_{2q}=\lambda_{2q}$ & $1\times10^{-2}$ & $5\times10^{-2}$ \\
$T_1$ (ns)                   & 100\,000 & 50\,000 \\
$T_2$ (ns)                   & 80\,000  & 30\,000 \\
$p_{\text{ro}}$              & $2\times10^{-2}$ & $5\times10^{-2}$ \\
\bottomrule
\end{tabular}%
}
\caption{Preset uniformi illustrativi per le campagne rumorose N15. Non rappresentano la calibrazione di una QPU reale.}
\label{tab:use_case_params}
\end{table}

\subsubsection{Razionale per la scelta delle istanze}
\label{subsec:razionale_istanze}

\textbf{N15 --- riferimento e stress.} La base $a=7$ ha periodo $r=4$ e consente di estrarre i fattori 3 e 5. Il moltiplicatore controllato implementa realmente $y\mapsto7y\bmod15$: l'operazione elementare viene ripetuta \texttt{power} volte, anziché selezionare erroneamente un ramo mediante \texttt{power \% 4}. Il circuito compilato con Qiskit~2.5, \texttt{optimization\_level=2}, seed di transpiler fissato e base \texttt{rz/sx/x/cx} ha profondità 412 e contiene 224 porte \texttt{cx}; questi dati e l'hash SHA-256 della netlist sono registrati nel manifest. Le rotazioni \texttt{rz} sono cambi di frame virtuali, con durata ed errore nulli; \texttt{sx} e \texttt{x} hanno durata 50\,ns e \texttt{cx} 300\,ns. Depolarizzazione e rilassamento termico sono composti soltanto su \texttt{sx/x} e \texttt{cx}; il readout è una matrice di confusione simmetrica.

\textbf{N21 e N35 --- validazione ideale.} Per $N=21$, $a=2$, e $N=35$, $a=6$, si usa l'aritmetica reversibile di Beauregard. La verifica controlla l'azione sui vettori di base validi, il ritorno a zero degli ancilla e la QPE ideale. Come audit di risorse, il circuito completo $N=21$, $n_{\text{count}}=8$, compilato globalmente nella stessa base a livello~2, conta 21\,036 \texttt{cx} e profondità 23\,081. Al valore illustrativo $\lambda=10^{-3}$, il proxy circuitale registrato è $2.6987\times10^{-9}$: indica accumulo circuitale, non la probabilità di successo di Shor. L'esclusione dal rumore è una scelta di validità e costo, non un risultato sulla robustezza delle due istanze.

% -----------------------------------------------
\subsection{Parametri di Input e Output}
\label{sec:input_output}

\subsubsection{Parametri di Input}

La Tabella~\ref{tab:input_params} elenca i parametri di configurazione del sistema. I loro valori vengono fissati per ciascun use case prima dell'esecuzione e rimangono costanti per tutta la durata dell'esperimento corrispondente.

\begin{table}[H]
\centering
\begin{tabular}{@{}lll@{}}
\toprule
\textbf{Parametro}  & \textbf{Unità} & \textbf{Descrizione} \\
\midrule
$N$                 & intero         & Numero da fattorizzare \\
$a$                 & intero         & Base della moltiplicazione modulare ($\gcd(a,N)=1$) \\
$n_{\text{count}}$  & intero         & Qubit del registro di controllo \\
$M_{\text{shots}}$  & intero         & Shot per esecuzione del simulatore \\
$\lambda_{1q}$      & adimensionale  & Parametro Aer del canale depolarizzante 1Q \\
$\lambda_{2q}$      & adimensionale  & Parametro Aer del canale depolarizzante 2Q \\
$T_1$               & ns             & Tempo di rilassamento (amplitude damping) \\
$T_2$               & ns             & Tempo di dephasing \\
$p_{\text{ro}}$     & adimensionale  & Probabilità di readout error \\
\bottomrule
\end{tabular}
\caption{Parametri principali del sistema sperimentale. $\lambda$ segue la convenzione Aer e non coincide con la probabilità aggregata di Pauli non identità.}
\label{tab:input_params}
\end{table}

\subsubsection{Output del Sistema}

Per ogni use case e per ciascuno dei due metodi, il sistema produce le seguenti quantità:

\begin{itemize}
    \item I fattori trovati $p$ e $q$ tali che $p \cdot q = N$, con verifica esplicita della correttezza ($p \neq 1$, $q \neq 1$, $p \cdot q = N$);
    \item Il numero di iterazioni necessarie per il primo risultato corretto;
    \item La distribuzione di frequenza completa degli output del registro di controllo su $M_{\text{shots}}$ esecuzioni;
    \item Le metriche di confronto tra i due metodi definite nella Sezione~\ref{sec:metriche}.
\end{itemize}

Per il solo Metodo 2, vengono prodotte in aggiunta:
\begin{itemize}
    \item L'accuratezza del classificatore e il punteggio F1 sul test set;
    \item La curva ROC con l'area sottostante (\ac{AUC}).
\end{itemize}

\subsection{Approfondimento della Sezione 2.6: Definizioni Complete delle Metriche}
\label{app:obiettivi:metriche}

\noindent\textit{Origine nel corpo:} definizioni operative estese della sintesi in Tabella~\ref{tab:contratto_sintesi} e della Sezione~\ref{subsec:metriche_qec}.

\subsubsection{Metriche della prima fase}

Per il \textbf{Metodo 1} si registrano: $\bar{M}_1$, numero medio di esecuzioni necessarie al primo risultato corretto sulle repliche riuscite, sempre accompagnato dal tasso di successo; $\sigma_{M_1}$, deviazione standard delle iterazioni; $P_{\mathrm{success},1}(M_{\max})$, probabilità di successo entro il budget preregistrato $M_{\max}=50$.

Per il \textbf{Metodo 2} si registrano analogamente $\bar{M}_2$, $\sigma_{M_2}$ e $P_{\mathrm{success},2}(M_{\max})$ sulle stesse repliche. La qualità predittiva è descritta da accuratezza, F1-score della classe positiva e \ac{AUC} sul test set indipendente. Tali metriche qualificano il classificatore, ma la metrica operativa del metodo resta il numero di iterazioni fino a una fattorizzazione verificata.

La metrica descrittiva principale è il fattore di riduzione $\rho=\bar{M}_1/\bar{M}_2$ dell'Equazione~\eqref{eq:fattore_riduzione}. $\rho>1$ indica soltanto una media inferiore per il metodo al denominatore: la significatività viene stabilita sulle differenze appaiate e non dedotta dal rapporto fra medie.

\subsubsection{Metriche QEC e analisi di sensibilità}

\begin{itemize}
    \item $p_L$: probabilità che il qubit logico subisca un errore non corretto dopo codifica, rumore, estrazione di sindrome e decodifica; è misurata in funzione dell'errore fisico $p$ e della distanza $d$.
    \item \textbf{Pendenza di soppressione}: esponente $\gamma$ nella relazione $p_L\propto p^\gamma$ a piccolo $p$; per distanza $d=3$ il comportamento atteso è $\gamma\simeq2$.
    \item \textbf{Break-even}: regione $p_L<p$, nella quale il qubit logico è più affidabile del qubit fisico non protetto.
    \item \textbf{Soglia $p_{th}$}: valore intorno al quale si incrociano le curve del surface code a distanze diverse; sotto soglia, aumentare $d$ deve ridurre $p_L$.
    \item $P_{\mathrm{success}}(p_L)$ e $P(k)$: in M8, $p_L$ denota la probabilità totale di un Pauli non identità dopo ogni gate compilato incluso nel proxy, non il tasso per ciclo dei benchmark QEC. La probabilità cumulativa dopo $k$ esecuzioni indipendenti è $P(k)=1-(1-P_s)^k$.
\end{itemize}

Ogni stima Monte Carlo riporta numero di campioni, seed e incertezza binomiale. La diversa definizione di $p_L$ nei memory experiment e in M8 viene sempre dichiarata per impedire confronti numerici non commensurabili.

\subsection{Approfondimento della Sezione 2.3: Protocollo Statistico e Test}
\label{app:obiettivi:protocollo}

\noindent\textit{Origine nel corpo:} dettagli dei test e dei criteri richiamati dopo le ipotesi e nella sintesi delle metriche del Capitolo~\ref{chap:obiettivi}.

\subsubsection{Appaiamento, fallimenti e pareggi}

Ogni replica genera un unico calendario di istogrammi: M1, $M_{\mathrm{TOP4}}$ e M2 analizzano gli stessi conteggi nello stesso ordine. Un fallimento entro $M_{\max}$ è codificato con la sentinella $M_{\max}+1$, così rimane distinto dal successo all'ultima iterazione e non scompare dai vettori inferenziali. Il confronto usa il test dei ranghi con segno di Wilcoxon con \texttt{zero\_method='pratt'} e soglia $\alpha=0.05$: M1$>$TOP-4 e M1$>$M2 sono unilaterali, TOP-4--M2 è bilaterale. Quando sono riportati più confronti della stessa famiglia, i valori $p$ vengono corretti con Holm.

In caso di conteggi uguali, l'ordine dei candidati non dipende né dall'ordine del dizionario né dal valore numerico della bitstring: è stabilito da una priorità SHA-256 dipendente dal seed della coppia replica--iterazione. Seed base, calendari distinti per simulatore e tie-breaking, configurazione e hash degli artefatti sono registrati nel manifest.

\subsubsection{Mitigazione e regole decisionali}

Il readout error è incorporato nel simulatore come matrice di confusione simmetrica parametrizzata da $p_{\mathrm{ro}}$; non viene applicata una correzione post-hoc. La sensibilità a $p_{\mathrm{ro}}$ è misurata nell'artefatto parametrico schema~2 corrente (Sezione~\ref{subsec:sweep_secondari}); gli artefatti precedenti alla correzione del circuito non alimentano i risultati.

M1 seleziona la sola moda dell'istogramma, senza sogliatura preliminare, e applica direttamente \texttt{extract\_factors}. TOP-4 estende il tentativo ai quattro esiti più frequenti. M2 usa l'etichetta TOP-1 per decidere se accettare l'istogramma e, in caso positivo, applica TOP-4: proprio il confronto con TOP-4 privo di filtro isola l'effetto del classificatore.

\subsubsection{Matrice di verifica}

La matrice sintetica che collega componenti, controlli e criteri di evidenza è
stata mantenuta nel corpo, nella Tabella~\ref{tab:verifica_requisiti}. Le
sottosezioni precedenti ne conservano qui le definizioni operative estese.

\end{onehalfspace}
